\ifdefined\iscombined
	\lecturetitle{Exercises in Understanding Design Basics}
\else
\documentclass{article}
\usepackage{listings}
\usepackage{amsthm}
\usepackage{comment}
\usepackage{amsmath}
\usepackage{amsfonts}
\usepackage{sectsty}
\usepackage{hyperref}
\usepackage{fancyhdr}
\usepackage[top=1in,bottom=1in,left=1in,right=1in]{geometry}
\usepackage{float}
\usepackage{graphicx}
\usepackage{tabularx}
\usepackage{mathtools}

\date{
	\vspace{-.75em}
	{\large \today}
}

\title{
	\vspace*{-1.5em}
	{\Huge Lecture 8} \\
	\vspace{-1em}
	\rule{\textwidth/4}{.01em} \\
	\vspace{-.25em}
	{\Large Exercises in Understanding Design Basics}
	\vspace{-.75em}
}

\author{
	{\large Matthew Farah}
}

\hypersetup{
	colorlinks=true,
	linkcolor=blue,
	filecolor=magenta,
	urlcolor=blue,
	citecolor=blue,
	anchorcolor=blue
}

\fancyhead{}
\fancyhead[L]{\today}
\fancyhead[R]{Lecture 8 - Exercises in Understanding Design Basics}
\sectionfont{\fontsize{12}{12}\bfseries}
\subsectionfont{\fontsize{10}{12}\normalfont}
\subsubsectionfont{\fontsize{10}{12}\normalfont}

\begin{document}

\maketitle
\pagestyle{fancy}

\fi

\emph{I didn't take notes for each example taken up in class for this one, just information that seemed generally relevant} \\

\begin{itemize}
	\item Choosing the form that a system should take is more challenging than what functionalities it should have.
	\item A system is a collection of one or more elements in an environment which fulfills one or more tasks.
	\begin{itemize}
		\item Composite systems are defined as systems which use more than one such element.
	\end{itemize}
	\item Modules are used to hide secrets.
	\item Business events are independent of one another.
\end{itemize}

\ifdefined\iscombined
	\newpage
\else
	\end{document}
\fi
